\documentclass[10pt,a4paper]{amsart}
\usepackage[margin=19mm]{geometry}
\usepackage{amsmath,amssymb,mathtools}
\usepackage[scr=rsfs,cal=euler]{mathalfa}
\usepackage{xfrac}
\usepackage[unicode,hidelinks,bookmarks=false]{hyperref}
\newcommand{\ie}{i.e.\ }
\newcommand{\cf}{cf.\ }
\newcommand{\resp}{respectively\ }
\newcommand{\Subsection}{\subsection}
\providecommand{\nobreakdash}{\nobreak\hbox{-}}
\setlength{\emergencystretch}{2em}
\multlinegap0pt
\usepackage{tikz}
\usepackage{amssymb}
\usepackage{mathtools}
\newtheorem{theorem}{Theorem}
\newcommand{\obsdata}{\train_{\text{obs}}}
\newcommand{\train}{\mathcal{D}}
\def\transpose{\intercal}
\title{The relative value of interventional and observational samples in Bayesian Causal Linear Gaussian Models}
\author{Valentinian Lungu, Anish Dhir, Mark van der Wilk, Ioannis Kontoyiannis}
\date{Source: arXiv:2603.26460v1 (CC BY 4.0)}
\begin{document}
\maketitle

\noindent\emph{Source and licence.} The relative value of interventional and observational samples in Bayesian Causal Linear Gaussian Models. Version arXiv:2603.26460v1 (CC BY 4.0), distributed by arXiv under the Creative Commons Attribution 4.0 International licence, http://creativecommons.org/licenses/by/4.0/, as recorded in the arXiv licence metadata for this exact version and shown on the abstract page https://arxiv.org/abs/2603.26460v1. Full licence notice: \texttt{licenses/arxiv-2603.26460-CC-BY-4.0.txt}.\par\smallskip

\noindent\textit{Editorial scope.} Counted unit: the article's theorem thm:BGe\_nonindent (source0.tex lines 1120--1136). The article's complete original proof of it is delivered with it (source0.tex lines 1137--1161, one \texttt{proof} environment of 25 lines). No mathematics is written, repaired or re-derived: the statement and the proof are the article's own lines, in the article's own notation, and no retained line is substituted or re-flowed.  Retained uncounted as the source-local context the proof needs: lines 130--133; lines 364--370; lines 682--685.  Nothing was omitted from any retained span, so the package carries no omission record.  No retained line contains a citation, so the wrapper carries no bibliography and no bibliography entry.  No retained line contains a figure environment, an image-inclusion command or drawing code, so no image is delivered and the package declares no asset dependency.  Editorial formatting only, in addition: an AMS \texttt{amsart} preamble that restates the article's own declarations and macros used by the retained lines, namely the environments \texttt{theorem} on separate counters, together with the standard packages those lines need.  The retained environments print in this document as Theorem 1 in order of appearance, whereas the article prints the same items with the numbering of its own full document.  The complete native source of the article as distributed in the arXiv e-print for this exact version is bundled unchanged as \texttt{sources/197284/source0.tex}.
\begin{equation}
    X = A^\transpose X + \epsilon,
    \label{eq:SCM_vector}
\end{equation}

\begin{equation}
\label{eq:BGe_priors_compact}
\begin{split}
    S \sim \text{Unif}\{S^1, S^2, S^3\}; \quad \tau_j^2 \mid S^i \sim \text{IG}(\alpha_{i,j}, \beta), \\
    w \mid (S^i, \tau^2) \sim N(0, \lambda\tau_i^2)^{\mathbb{I}(i \neq 3)}
\end{split}
\end{equation}

    \begin{equation}
        S^X_{12} := \sum_{i=1}^n X_i(1)X_i(2), \quad S^X_1:= \sum_{i=1}^n X_i(1)^2, \quad S^X_2:= \sum_{i=1}^n X_i(2)^2,
        \label{eq:MLE_interv_X}
    \end{equation}

\begin{theorem}
    \label{thm:BGe_nonindent}
    Let $\obsdata = (X_i)_{1\leq i\leq n}$ be i.i.d. observations generated by \eqref{eq:SCM_vector}. Given the hierarchical priors defined in \eqref{eq:BGe_priors_compact}, the posterior distribution over the structures, represented by the vector $[\pi(S^1 \mid \obsdata), \pi(S^2 \mid \obsdata), \pi(S^3 \mid \obsdata)]^\transpose$, is
    \begin{align*}
        \Bigg[ 
        & \frac{(2\beta)^{\alpha_1+\alpha_2} \Gamma(\alpha_2+\frac{n}{2})\Gamma(\alpha_1+\frac{n}{2})}{\sqrt{\eta} \pi^n \Gamma(\alpha_1)\Gamma(\alpha_2)} \cdot \frac{(S_2^{X\lambda})^{\alpha_1+(n-1)/2}}{(S_2^{X\beta})^{\alpha_2+n/2}} \cdot \frac{1}{(S_1^{X\beta} S_2^{X\lambda} - (S_{12}^X)^2)^{\alpha_1+n/2}}, \\
        & \frac{(2\beta)^{\alpha_3+\alpha_4} \Gamma(\alpha_3+\frac{n}{2})\Gamma(\alpha_4+\frac{n}{2})}{\sqrt{\eta} \pi^n \Gamma(\alpha_3)\Gamma(\alpha_4)} \cdot \frac{(S_1^{X\lambda})^{\alpha_4+(n-1)/2}}{(S_1^{X\beta})^{\alpha_3+n/2}} \cdot \frac{1}{(S_1^{X\lambda} S_2^{X\beta} - (S_{12}^X)^2)^{\alpha_4+n/2}}, \\
        & \frac{(2\beta)^{\alpha_5+\alpha_6} \Gamma(\alpha_5+\frac{n}{2})\Gamma(\alpha_6+\frac{n}{2})}{\pi^n \Gamma(\alpha_5)\Gamma(\alpha_6)} \cdot \frac{1}{(S_1^{X\beta})^{\alpha_5+n/2}(S_2^{X\beta})^{\alpha_6+n/2}} 
        \Bigg]^\transpose,
    \end{align*}
    where 
    \begin{equation}
        S_i^{X\lambda} := S_{i}^X + \frac{1}{\lambda}, \quad S_{i}^{X\beta} := S_{i}^X + 2\beta \quad \text{for } i \in \{1, 2\},
        \label{eq:def_S_beta_eta}
    \end{equation}
    and $S_1^X, S_2^X$, and $S_{12}^X$ are defined in \eqref{eq:MLE_interv_X}.
\end{theorem}

\begin{proof}
    \begin{align}
        \notag \pi(S^1\mid\obsdata) &\propto \int f(\obsdata\mid S^1, w, \tau_1^2, \tau_2^2) f(w\mid\tau_1^2)f(\tau_1^2)f(\tau_2^2)\, d(\tau_1^2)\, d(\tau_2^2)\, dw \\
        \notag &= \int \frac{1}{(\sqrt{2\pi \tau_1^2})^n}\exp \left(-\frac{1}{2\tau_1^2}\sum_{i=1}^n(X_1^{i}-wX_2^{i})^2 \right) \frac{1}{(\sqrt{2\pi \tau_2^2})^n}\exp \left(-\frac{1}{2\tau_2^2}S_2^X \right) \\
        \notag &\quad \times \frac{1}{\sqrt{2\pi\lambda\tau_1^2}}\exp\left(-\frac{w^2}{2\lambda\tau_1^2} \right) \frac{\beta^{\alpha_1+\alpha_2}}{\Gamma(\alpha_1)\Gamma(\alpha_2)} (\tau_1^2)^{-\alpha_1-1} \exp \left(-\frac{\beta}{\tau_1^2} \right) (\tau_2^2)^{-\alpha_2-1} \exp \left(-\frac{\beta}{\tau_2^2} \right)\, d(\tau_1^2)\, d(\tau_2^2)\, dw \\
        \notag &= \frac{1}{\sqrt{\lambda}} \frac{1}{(2\pi)^{n+1/2}} \frac{\beta^{\alpha_1+\alpha_2}}{\Gamma(\alpha_1)\Gamma(\alpha_2)} \frac{\Gamma(\alpha_2+\frac{n}{2})}{\left( \beta + \frac{S_2^X}{2} \right)^{\alpha_2+n/2}}\times \int_{-\infty}^\infty \frac{\Gamma(\alpha_1+\frac{n+1}{2})}{\left( \beta+ \frac{\sum_{i=1}^n(X_1^{i}-wX_2^{i})^2}{2} +\frac{w^2}{2\lambda}   \right)^{\alpha_1+\frac{n+1}{2}}} \, dw \\
        \notag &= \frac{(2\beta)^{\alpha_1+\alpha_2} \Gamma(\alpha_2+\frac{n}{2})\Gamma(\alpha_1+\frac{n}{2})}{\sqrt{\lambda} \pi^n \Gamma(\alpha_1)\Gamma(\alpha_2)} \cdot \frac{(1/\lambda + S_2^X)^{\alpha_1+(n-1)/2}}{(2\beta + S_2^X)^{\alpha_2+n/2}} \cdot \frac{1}{[(2\beta+S_1^X)(1/\lambda+S_2^X) - (S_{12}^X)^2]^{\alpha_1+n/2}}\\
        &= \frac{(2\beta)^{\alpha_1+\alpha_2} \Gamma(\alpha_2+\frac{n}{2})\Gamma(\alpha_1+\frac{n}{2})}{\sqrt{\lambda} \pi^n \Gamma(\alpha_1)\Gamma(\alpha_2)} \cdot \frac{ (S_2^{X\lambda})^{\alpha_1+(n-1)/2}}{(S_2^{X\beta})^{\alpha_2+n/2}} \cdot \frac{1}{(S_1^{X\beta} S_2^{X\lambda} - (S_{12}^X)^2)^{\alpha_1+n/2}},
        \label{eq:post_G1}
    \end{align}

    Similarly, we have
    \begin{equation}
        \pi(S^2\mid\obsdata) = \frac{(2\beta)^{\alpha_3+\alpha_4} \Gamma(\alpha_3+\frac{n}{2})\Gamma(\alpha_4+\frac{n}{2})}{\sqrt{\lambda} \pi^n \Gamma(\alpha_3)\Gamma(\alpha_4)} \cdot \frac{ (S_1^{X\lambda})^{\alpha_4+(n-1)/2}}{(S_1^{X\beta})^{\alpha_3+n/2}} \cdot \frac{1}{(S_1^{X\lambda} S_2^{X\beta} - (S_{12}^X)^2)^{\alpha_4+n/2}},
        \label{eq:post_G2}
    \end{equation}
    and
    \begin{align}
        \notag \pi(S^3\mid\obsdata) &\propto \int f(\obsdata\mid S^3, \tau_1^2, \tau_2^2) f(\tau_1^2)f(\tau_2^2)\, d(\tau_1^2)\, d(\tau_2^2) \\
        \notag &= \int \frac{1}{(\sqrt{2\pi \tau_1^2})^n}\exp \left(-\frac{1}{2\tau_1^2}S_1^X \right) \frac{1}{(\sqrt{2\pi \tau_2^2})^n}\exp \left(-\frac{1}{2\tau_2^2}S_2^X \right)\\
        \notag &\quad \times \frac{\beta^{\alpha_5+\alpha_6}}{\Gamma(\alpha_5)\Gamma(\alpha_6)} (\tau_1^2)^{-\alpha_5-1} \exp \left(-\frac{\beta}{\tau_1^2} \right) (\tau_2^2)^{-\alpha_6-1} \exp \left(-\frac{\beta}{\tau_2^2} \right)\, d(\tau_1^2)\, d(\tau_2^2)\\
        &= \frac{(2\beta)^{\alpha_5+\alpha_6} \Gamma(\alpha_5+\frac{n}{2})\Gamma(\alpha_6+\frac{n}{2})}{ \pi^n \Gamma(\alpha_5)\Gamma(\alpha_6)} \cdot \frac{1}{ (S_1^{X\beta})^{\alpha_5+n/2}(S_2^{X\beta})^{\alpha_6+n/2}}.
        \label{eq:post_G3}
    \end{align}
\end{proof}

\end{document}
